% Options for packages loaded elsewhere
\PassOptionsToPackage{unicode}{hyperref}
\PassOptionsToPackage{hyphens}{url}
\documentclass[
  english,
]{article}
\usepackage{xcolor}
\usepackage{amsmath,amssymb}
\setcounter{secnumdepth}{5}
\usepackage{iftex}
\ifPDFTeX
  \usepackage[T1]{fontenc}
  \usepackage[utf8]{inputenc}
  \usepackage{textcomp} % provide euro and other symbols
\else % if luatex or xetex
  \usepackage{unicode-math} % this also loads fontspec
  \defaultfontfeatures{Scale=MatchLowercase}
  \defaultfontfeatures[\rmfamily]{Ligatures=TeX,Scale=1}
\fi
\usepackage{lmodern}
\ifPDFTeX\else
  % xetex/luatex font selection
\fi
% Use upquote if available, for straight quotes in verbatim environments
\IfFileExists{upquote.sty}{\usepackage{upquote}}{}
\IfFileExists{microtype.sty}{% use microtype if available
  \usepackage[]{microtype}
  \UseMicrotypeSet[protrusion]{basicmath} % disable protrusion for tt fonts
}{}
\makeatletter
\@ifundefined{KOMAClassName}{% if non-KOMA class
  \IfFileExists{parskip.sty}{%
    \usepackage{parskip}
  }{% else
    \setlength{\parindent}{0pt}
    \setlength{\parskip}{6pt plus 2pt minus 1pt}}
}{% if KOMA class
  \KOMAoptions{parskip=half}}
\makeatother
\usepackage{longtable,booktabs,array}
\usepackage{calc} % for calculating minipage widths
% Correct order of tables after \paragraph or \subparagraph
\usepackage{etoolbox}
\makeatletter
\patchcmd\longtable{\par}{\if@noskipsec\mbox{}\fi\par}{}{}
\makeatother
% Allow footnotes in longtable head/foot
\IfFileExists{footnotehyper.sty}{\usepackage{footnotehyper}}{\usepackage{footnote}}
\makesavenoteenv{longtable}
\ifLuaTeX
\usepackage[bidi=basic,shorthands=off,]{babel}
\else
\usepackage[bidi=default,shorthands=off,]{babel}
\fi
\ifLuaTeX
  \usepackage{selnolig} % disable illegal ligatures
\fi
\setlength{\emergencystretch}{3em} % prevent overfull lines
\providecommand{\tightlist}{%
  \setlength{\itemsep}{0pt}\setlength{\parskip}{0pt}}
\usepackage{bookmark}
\IfFileExists{xurl.sty}{\usepackage{xurl}}{} % add URL line breaks if available
\urlstyle{same}
\hypersetup{
  pdftitle={Incarnation and Unitarity},
  pdfauthor={Paul Koop},
  pdflang={en},
  hidelinks,
  pdfcreator={LaTeX via pandoc}}

\title{Incarnation and Unitarity}
\usepackage{etoolbox}
\makeatletter
\providecommand{\subtitle}[1]{% add subtitle to \maketitle
  \apptocmd{\@title}{\par {\large #1 \par}}{}{}
}
\makeatother
\subtitle{Christology after the End of Theodicy -- From the Limit Value
to the Present: Sacrament, Church and Diakonia in Evolutionary Reality}
\author{Paul Koop}
\date{2026}

\begin{document}
\maketitle

{
\setcounter{tocdepth}{3}
\tableofcontents
}
\section{Incarnation and Unitarity}\label{incarnation-and-unitarity}

\subsection{Christology after the End of Theodicy -- From the Limit
Value to the
Present}\label{christology-after-the-end-of-theodicy-from-the-limit-value-to-the-present}

\subsubsection{Sacrament, Church and Diakonia in Evolutionary
Reality}\label{sacrament-church-and-diakonia-in-evolutionary-reality}

\begin{center}\rule{0.5\linewidth}{0.5pt}\end{center}

\textbf{Table of Contents}

\begin{itemize}
\tightlist
\item
  \hyperref[-1-finitude-and-being-the-fundamental-philosophical-question]{§
  1. Finitude and Being -- The Fundamental Philosophical Question}
\item
  \hyperref[-2-epistemology-of-the-search-for-the-limit-value-how-do-we-know]{§
  2. Epistemology of the Search for the Limit Value -- How Do We Know?}
\item
  \hyperref[-3-nothingness-as-creative-limit-value-from-what-do-we-know]{§
  3. Nothingness as Creative Limit Value -- From What Do We Know?}
\item
  \hyperref[-4-the-cosmological-trinity]{§ 4. The Cosmological Trinity}
\item
  \hyperref[-5-incarnation-the-historical-form-of-undivided-reality]{§
  5. Incarnation -- The Historical Form of Undivided Reality}
\item
  \hyperref[-6-cross-and-suffering-unitarity-of-pain-and-the-end-of-theodicy]{§
  6. Cross and Suffering -- Unitarity of Pain and the End of Theodicy}
\item
  \hyperref[interlude-song-of-the-multiverse]{Interlude: Song of the
  Multiverse}
\item
  \hyperref[-7-resurrection-and-omega-point-hope-without-metaphysical-shortcut]{§
  7. Resurrection and Omega Point -- Hope without Metaphysical Shortcut}
\item
  \hyperref[-8-sacrament-and-eucharist-presence-transformation-material-reality]{§
  8. Sacrament and Eucharist -- Presence, Transformation, Material
  Reality}
\item
  \hyperref[-9-church-and-liturgy-community-without-possession-of-the-limit-value]{§
  9. Church and Liturgy -- Community without Possession of the Limit
  Value}
\item
  \hyperref[-10-diakonia-and-social-ethics-freedom-responsibility-reduction-of-suffering]{§
  10. Diakonia and Social Ethics -- Freedom, Responsibility, Reduction
  of Suffering}
\item
  \hyperref[-11-diagnosis-of-the-present-ai-secularity-technology-pastoral-care]{§
  11. Diagnosis of the Present -- AI, Secularity, Technology, Pastoral
  Care}
\item
  \hyperref[notes]{Notes}
\item
  \hyperref[glossary]{Glossary}
\item
  \hyperref[summary]{Summary}
\item
  \hyperref[editorial-note]{Editorial Note}
\end{itemize}

\begin{center}\rule{0.5\linewidth}{0.5pt}\end{center}

\subsection{§ 1. Finitude and Being -- The Fundamental Philosophical
Question}\label{finitude-and-being-the-fundamental-philosophical-question}

Western metaphysics begins with a distinction that it has never entirely
overcome: the distinction between finite, mutable being and infinite,
absolute being. Finite being was considered deficient, permeated by
non-being, constantly threatened by annihilation. The creature's longing
was directed toward escape from temporality into eternal immutability.
Augustine translated this longing into the language of inwardness: God
is more inward than my innermost being, and at the same time the
absolutely sublime. Thomas Aquinas translated it into the language of
substance: God as pure being, which knows no change, no suffering, no
finitude.

This perspective, however, misjudges the ontological dignity of
finitude. Finitude is not a lack of being, but the only possible form in
which personal freedom, the capacity for relationship, and concrete
history can exist. Without temporality there is no action, without
limitation no decision, without vulnerability no love. The redefinition
of being must therefore proceed from finitude as a positive ontological
category.

Anyone who takes finitude seriously must also rethink the concept of
freedom. Freedom must not be misunderstood as the suspended, unbound
arbitrariness of an isolated subject. It is always situated freedom,
bound to the bodily, biological, and historical constitution of the
human being. We call this structure \textbf{Gaussian freedom}: personal
freedom of decision does not unfold in a value-free, infinite space, but
within the universal normal distribution and the material limits of the
cosmos. It is the capacity of the person to write a unique, indelible
life story within a reality that exhibits probabilistic and
deterministic structures. Freedom is not the absence of causes, but the
absence of disturbance. It is the undisturbed execution of one's own
curve.

Being is thus not to be understood as a resting substrate layer beneath
things, but as the network of relationships itself, in which entities
arise, develop, and interact with one another. When being is understood
relationally, the question of the ultimate ground of reality also
changes. God is no longer the infinitely distant super-object, but the
sustaining depth and the absolute limit toward which every finite
relationship is designed.

Here the \textbf{Eastern doctrine of theosis} becomes fruitful. It
speaks of transfiguration, of unity, of the healing dimension of the
divine light. Deification is not the overcoming of finitude, but its
fulfillment. The human being does not become God by ceasing to be human,
but by becoming fully human -- open to the limit value that sustains
him. The Western question of guilt and grace and the Eastern question of
transfiguration and unity are not opposites, but complementarities.
Augustine took the inner conflict seriously; the Eastern tradition took
transfiguration seriously. Both saw something that the other tradition
overlooked. A Christology after evolution must integrate both: the
reality of guilt and the possibility of transfiguration.

The fundamental philosophical question concerning the connection between
finitude and being thus leads directly into the center of an
evolutionary Christology. If being itself is structured relationally and
processually, the absolute principle cannot be thought in isolation from
the world. Theology must leave behind the aporias of substance thinking
and develop an ontology in which the interweaving of creator and
creation becomes conceivable without extinguishing the autonomy of the
creature. This ontology is the ontology of the \textbf{limit value}: not
a distant goal, but a reality that is already effective in every
approximation.

The Augustinian tradition is thereby appreciated without being followed.
Augustine took seriously the experience of inner conflict, the need for
grace, inwardness -- that remains. But his deficit ontology, his
original-sin punishment, his notion of an angry judge who saves a few
through grace -- that is rejected. Not out of polemic, but because it
does not do justice to finitude. Finitude is not a curse, but the
condition of the possibility of freedom, love, and history. Whoever
understands it as a curse understands neither finitude nor redemption.

\begin{center}\rule{0.5\linewidth}{0.5pt}\end{center}

\subsection{§ 2. Epistemology of the Search for the Limit Value -- How
Do We
Know?}\label{epistemology-of-the-search-for-the-limit-value-how-do-we-know}

The recognition of evolution as a universal principle of reality compels
not only theology but also epistemology to revise its presuppositions.
If reality is not static but historical, processual, and open, then
knowledge too cannot be understood as the possession of truth, but only
as approximation. This approximation is asymptotic: it approaches the
limit value without ever fully reaching it.

The starting point is the assumption of an objective reality that exists
independently of the knowing subject. This assumption is not provable,
but pragmatically indispensable: without it, criticism, error, and
progress would be meaningless. At the same time, the knowing subject is
an integral part of reality. Knowledge is a natural, causally embedded
process that overcomes Cartesian dualism. Consciousness is not the
origin of reality, but a product of real processes. It arises within the
world, not above it.

Ontologically real is only the present. Past and future do not exist as
things, but as epistemic horizons: the past is reconstructed, never
fully accessible; the future is anticipated, in principle open. In the
present, causal necessity and open possibility overlap. Decision and
freedom do not arise through a break in causality, but as
self-organization within this overlap. Freedom is not a metaphysical
state of exception, but a highly complex form of intra-worldly order.

Past and future can be understood as limit values that knowledge
asymptotically approaches without ever fully reaching them. These limit
values are: the \textbf{origin} (the epistemically never fully
reconstructible beginning of all reality), the \textbf{present} (the
only place of real actualization), and the \textbf{Omega} (the conceived
endpoint of complete coherence and self-transparency of reality). These
three limit values are not additional entities, but perspectives on one
and the same reality. Knowledge moves between them without ever
possessing them.

Truth in this model is not timelessly given, but historically proven. It
arises through error, correction, and the stabilization of viable
patterns of interpretation. This epistemology thus stands close to a
fallibilist concept of truth. Absolute certainty is excluded, but
increasing truth content is possible. This structure allows revelation
to be understood not as a break in the order of knowledge, but as an
event of heightened coherence that must prove itself historically.

Here the \textbf{argumentative engagement with critical rationalism and
the many-worlds interpretation} is indispensable. Critical rationalism
provides the epistemological foundation: knowledge is always
provisional; approximation occurs through falsification; the open
society is the political form of this fallibility. The many-worlds
interpretation provides the ontological foundation: possibilities are
real, but not separate and substantial; reality is a network of
branchings in which each version has its own perspective. This work
agrees with critical rationalism and the many-worlds interpretation in
the rejection of dualism, in the emphasis on fallibility, in the high
estimation of the open society.

But it contradicts them at one decisive point: the infinite progress
postulated by critical rationalism is not the last word. Progress is
infinite -- but it is not aimless. It moves toward a limit value that is
already effective in every approximation. The many-worlds interpretation
provides the language, not the goal. It provides the method, not the
structure. This work takes it at its word and goes one step further than
it itself goes: it thinks the limit value as a boundary, not as a goal.

The plausibility check against physics is important, but it remains
analogy, not proof. Thermodynamics and entropy: local order (life,
consciousness) arises within dissipative structures, despite global
entropy increase. The limit value can be interpreted as a metaphysical
maximum order without violating the physical framework. Conservation of
energy: cognitive processes are embedded in the existing energy of the
universe; the conservation laws are respected. Quantum mechanics:
superposition and branching in physics serve as an analogy to the
epistemological model of the present, not as an exact physical homology.
Black-hole physics: information conservation at the event horizon is the
physical image of the unitary preservation of personal information.

\textbf{Physics = formal model, theology = content-related
interpretation.} This distinction is methodologically reflected. Physics
provides no proofs for theological statements; it provides structures
that can be interpreted theologically. Whoever blurs this distinction
loses both physics and theology.

\begin{center}\rule{0.5\linewidth}{0.5pt}\end{center}

\subsection{§ 3. Nothingness as Creative Limit Value -- From What Do We
Know?}\label{nothingness-as-creative-limit-value-from-what-do-we-know}

Nothingness is not the absence of existence. It is a boundary condition
that defines both the origin and the goal of reality. It is absolutely
homogeneous and does not differ from itself -- and precisely thereby it
is the foundation of everything that exists. This determination of
nothingness is neither nihilistic nor Neoplatonic. It would be
nihilistic if nothingness were annihilation, the final extinction of all
information. It would be Neoplatonic if nothingness were the ground of
emanation from which being pours forth in stages. Neither is meant.
Nothingness is the \textbf{limit value of self-concretization}: the
point at which possibilities condense into reality without an external
creator needing to intervene.

In the beginning, nothingness is the origin of all causal chains. From
it, reality arises through self-reference and self-relation. In the end,
everything returns to nothingness, since no differentiation exists that
goes beyond this limit value. In the present, nothingness becomes
concrete through the decoherence of countless possibilities and
manifests itself in consciousness as lived reality. Nothingness is not
the other of God; it is the manner in which God concretizes himself as
limit value.

The relationship to the concept of information is decisive. Information
is not nothingness, but it springs from it as limit value. Preservation
of information is the manner in which nothingness concretizes itself. If
the universe is unitary, then nothing may be finally lost. Neither a
lived thought nor a suffered sorrow -- and not the unlived possibilities
either. For they too are information. Every decision excludes
alternatives, but it does not erase them. The unrealized paths are not
trivialities, but real branchings that become invisible from the
perspective of one version. Their information remains part of the
overall structure.

Decoherence is the transition from the nothingness of possibilities to
the presence of experience. In quantum mechanics, decoherence designates
the process by which a quantum-mechanical system loses its superposition
and passes into a defined state. Theologically translated: decoherence
is the manner in which nothingness condenses into presence. The present
is the only ontologically real moment in which causal necessity and open
possibility overlap. Consciousness is the ongoing process of internal
pattern reinforcement within this overlap. The coherence of the I arises
from the stability of recurring patterns, without any need for a
reduction to neural mechanics or a metaphysical soul.

Nothingness as creative limit value is the ontological foundation of
unitarity. It explains why reality is closed without being static; why
it is open without being arbitrary; why it generates information without
losing it. It is the point at which language reaches its limit -- and at
which poetry begins.

\begin{quote}
In the beginning was possibility, infinitely deep, infinitely wide. A
whisper forms the first light, and life grows -- it contradicts the
rigidity of empty space, it winds itself from time and dream, it mirrors
itself, recognizes itself, consciousness that holds the boundaries.
\end{quote}

\begin{center}\rule{0.5\linewidth}{0.5pt}\end{center}

\subsection{§ 4. The Cosmological
Trinity}\label{the-cosmological-trinity}

The Trinity is not a dogma that one must believe without understanding
it. It is the \textbf{inner relationing of the one divine life}, which
unfolds in beginning, present, and end without dividing itself. Father =
almighty beginning, Son = omniscient end, Holy Spirit = omnipresent
present. This assignment is not modalistic, because it does not make the
three persons mere manifestations of a single God. It is also not
unitary in the historical sense, because it does not deny the divinity
of Christ. It is \textbf{unitary in the ontological sense}: God and
world are not two competing substances, but the entire creation
possesses its medium of existence in God, without pantheistically
merging with his essence.

The Father is the origin of all possibilities. He is not a creator in
the dualistic sense, who would have had the possibility of deciding
against creation. He is the origin from which everything proceeds,
without his having to intervene arbitrarily. The Son is the omniscient
end. He is not a distant judge, but the completion in which all
information is integrated. The Spirit is the present, in which
nothingness concretizes itself and possibilities condense into
experience.

\textbf{The Spirit as present} is the key to understanding the Trinity
after evolution. The Spirit is not an additional entity that intervenes
in the world from outside. He is the present itself -- the manner in
which God is present in the now. In the present, past and future
overlap; in the present, possibilities condense into reality; in the
present, nothingness becomes concrete. The Spirit is the life-giving
power that works in all possible worlds. He is not bound to space and
time, but he works in space and time. He is of one essence with the
beginning and the end, and yet he differs from both -- not as a separate
substance, but as the manner in which unity makes itself present.

The demarcation against \textbf{Unitarianism} (historical) is important:
historical Unitarianism was a rationalistic, anti-Trinitarian movement
that reduced the divine to a monolithic, impersonal one God and rejected
the divinity of Christ. Our unitarity, by contrast, is a strict
ontological nondualism: it states that God and world are not two
competing substances, but that the entire creation possesses its medium
of existence in God, without pantheistically merging with his essence.
The demarcation against \textbf{modalism} is equally important: modalism
teaches that Father, Son, and Spirit are only manifestations of a single
God. The cosmological Trinity, by contrast, teaches that the three
persons are truly distinguished -- not as separate substances, but as
different modes of relationing the one divine life.

The Trinity is the precondition for God to be rich in relationship
within himself and thereby able to receive the world into himself. A God
without inner relation would not be capable of creating relationship. A
God without difference would not be capable of bearing difference. The
Trinity is the ontological condition of the possibility of creation,
incarnation, and consummation.

\begin{center}\rule{0.5\linewidth}{0.5pt}\end{center}

\subsection{§ 5. Incarnation -- The Historical Form of Undivided
Reality}\label{incarnation-the-historical-form-of-undivided-reality}

Incarnation must be misunderstood neither in the sense of a mythological
descent of a heavenly being nor in the sense of a merely moral metaphor.
It is the \textbf{historical form} in which the possibility of an
undivided relationship between God, world, and finite personal reality
shows itself without abolishing the finitude of the world. Unitarity
designates the closed, information-theoretical, and ontological
wholeness of reality. In a unitary universe, no information, no lived
moment, and no suffered pain is lost.

\textbf{Jesus of Nazareth} is the historical form in which the limit
value is incarnate. He is not the metaphysical starting point of the
world, but its historical culmination point. In his life, his loving,
and his suffering it becomes manifest that God does not isolate himself
against matter, but that the divine life manifests itself precisely
there where a human being opens himself radically for the Thou and the
world. Jesus is the embodied \textbf{conscientia} -- the co-knowledge of
God and cosmos. In him the \textbf{incarnational continuity} becomes
tangible: the historical Jesus of Nazareth is not an isolated miracle,
but the culminating breakthrough of the cosmic movement in which matter
becomes capable of responding to God in perfect personal freedom.

His \textbf{freedom} is not the arbitrariness of an isolated subject,
but the undisturbed unfolding of his curve within the divine normal
distribution. He does what he must do because he is entirely himself.
His \textbf{love} is not a feeling, but the radical assent to the
unavailability of history. He loves by opening himself to the Thou
without possessing it. His \textbf{suffering} is not a sacrifice that
reconciles God, but the consequence of his uncompromising love in a
world dominated by violence. His \textbf{death} is not the end, but the
transition from visibility into complete integration. His
\textbf{resurrection} is not the return to a biological life, but the
confirmation of unitary preservation: what was once real cannot
disappear.

\textbf{How ontology and Christology go together}: unitary ontology
provides the structure: limit value, information, unitarity, Gaussian
freedom, unitarity of pain. Inverse Christology provides the form:
emergent interpretive form, incarnational continuity, sacramentality as
real continuation of the incarnation, social ethics as incarnation test.
This work combines both: \textbf{Christ as the historical form of the
limit value}. That is: the limit value is not an abstract idea
subsequently applied to Jesus. Jesus is the form in which the limit
value becomes historical. In him it shows what the limit value is: not a
distant goal, but a reality that is already effective in every
approximation. The incarnation is the historical form of undivided
reality -- not because Jesus would be a special case, but because he is
the case in which reality becomes transparent to itself.

This Christology contradicts neither the Council of Chalcedon nor the
doctrine of the preexistence of the Logos. It does, however, shift their
epistemological location. The unity of divine and human nature is not
deduced from above, but inferred from below. Ontological statements
follow historical experience, not the other way around. The preexistence
of Christ is not a physical before, but a relational statement: God
relates from eternity to Christ. It is a statement about the
faithfulness of God, not about a temporal prior state.

The incarnation is not a singular event, but a principle. Where reality
is accepted without being shortened, incarnation happens. Jesus is not
the exception to this rule, but its full form. In him it becomes
manifest that matter is not the opposite pole to spirit, but the
physical form in which love can give itself away.

\begin{center}\rule{0.5\linewidth}{0.5pt}\end{center}

\subsection{§ 6. Cross and Suffering -- Unitarity of Pain and the End of
Theodicy}\label{cross-and-suffering-unitarity-of-pain-and-the-end-of-theodicy}

Suffering in the world is the millstone on which every classical
theodicy is crushed. Earthquakes, cancer, the cruelty of evolution, in
which the survival of one is the death sentence of another, and finally
the abysses of human guilt -- all this cannot be fitted into the schema
of a ``kind educator.'' The question: ``How can God permit this?'' is
wrongly posed, because it proceeds from a God who could intervene from
outside in a mechanical system without suspending the laws of finitude.
Finitude necessarily implies vulnerability, friction, and loss. An
evolutionary system that is to bring forth complex living beings cannot
exist without stochasticity, without mutation, without aging, and
without pain. Pain is the biological reverse side of sentience; grief is
the ontological reverse side of love.

We therefore introduce the concept of the \textbf{unitarity of pain}.
Suffering is not a moral scandal that refutes the existence of God, and
not a punishment for sin, but the structural phenomenon of friction in
the unfolding of the evolutionary process. Suffering is information. It
marks places of maximum tension, maximum condensation, maximum
difference. Where nothing is at stake, nothing relevant happens. Where
suffering arises, states are decided, structures are stabilized or
discarded, paths are fixed. Suffering is the price of the system not
standing still. The more complex the structure, the greater the capacity
for suffering. Not as punishment, but as indicator.

\textbf{On the cross of Jesus of Nazareth} this unitarity of pain is
theologically sealed. The cross is not the place where an enraged God
had to be appeased by the blood of his son. Such a notion makes God into
a monster. The cross is the historical consequence of Jesus'
uncompromising love in a world dominated by violence. By God dying in
the crucified Jesus, he himself enters into the darkest extreme of pain,
of Godforsakenness, and of annihilation. The cross is the radical assent
to the unavailability of history: God does not save the world from its
contingency, but in it. Anything else would be mythology.

The \textbf{wounds of the Risen One} are the eternal proof that history
with all its horrors is preserved in God. Not defused, not corrected,
but fully remembered. The pain remains as real information in the heart
of the Trinity. The wounds are not a symbol of suffering, but suffering
itself -- in its complete integration. What was once real cannot
disappear. Not the pain either. Not the wound either. It remains, but it
remains as memory, not as torment.

\textbf{The end of theodicy} does not mean becoming indifferent to
suffering. It means, on the contrary, no longer explaining suffering
away. There is no higher meaning in the suffering of a child; there is
only the radical obligation to reduce suffering and the trust that even
the most senseless dying is not lost in the unitary wholeness of God.
\textbf{Reduction of suffering is the highest principle.} Not
justification, not excuse, not trivialization -- but the practical work
of reducing suffering. Every glass of water handed to the thirsty, and
every political structure that promotes justice, is a direct building up
of the body of Christ.

Theodicy fails because it understands suffering as an argument against
meaning. In fact, suffering is the price of meaning. A completely
suffering-free system would be a completely undifferentiated system. It
would have no inner perspective, no history, no individuality. It would
be still -- not redeemed, but empty. God as the limit value of the
process is not a consolation in the emotional sense. He promises no
sparing, no exception, no final justice in the moral register. He is
that toward which knowledge moves asymptotically without ever fully
reaching it -- but already effective in the approximation. Everything
that happens belongs to this path, precisely because nothing can be
omitted. Not despite suffering, but including suffering.

Therewith theodicy ends. Not because suffering would be explained or
excused, but because the question itself loses its hold. There is no
judge outside the system to whom one could present suffering. There is
only the process itself -- and its asymptotic approximation to the limit
value. God is not the one who should have prevented suffering. God is
that which exists only because nothing was prevented -- and which, as
the limit value of all knowledge at the end of all paths, proves to be
that which always already was.

\begin{center}\rule{0.5\linewidth}{0.5pt}\end{center}

\subsection{Interlude: Song of the
Multiverse}\label{interlude-song-of-the-multiverse}

\begin{quote}
In the beginning was not space, not time, only the possible in darkness.
A whisper broke through the seams, where no reality yet struggled.

Then rose the wave, gentle and far, melting the ground of solitude,
divided, broke -- ever anew, till being scattered into a thousand
worlds.

And every reality split in two, where light bent, where time inclined,
and with it too, in quiet sense, the question of origin.

Omega point, a spirit of light, woven from number, pure and plain,
measured all being in clear sight, yet saw not what withdrew.

For measurement broke the silent path, let worlds emerge from the
shadows, and what it knew vanished in the light -- knowledge grew, yet
it did not see.

Then the multiverse closed wide into the play of eternity, where spirit
sang from many voices, and every truth found clarity.

No knowledge died, no light faded, it fans out in a wide ring, and life
that recognizes itself returns home to the net that guides all.

A point as great as eternity, where multiplicity thrives within itself,
where spirit and world rest in clarity -- the last step to being as
doing.

And from the depth softly shines a new song -- from dark night, in a
thousand voices it sounds pure: The goal of being -- is to be one.
\end{quote}

\begin{center}\rule{0.5\linewidth}{0.5pt}\end{center}

\subsection{§ 7. Resurrection and Omega Point -- Hope without
Metaphysical
Shortcut}\label{resurrection-and-omega-point-hope-without-metaphysical-shortcut}

How can the resurrection of the dead be thought under the conditions of
a strictly unitary, information-preserving ontology? The classical
notion of a biological reanimation or an otherworldly survival falls
short. It is the \textbf{constitutional duty of a unitary universe to
preserve personal information}.

The logic of information preservation is the key. If in a unitary cosmos
no information can be destroyed -- analogous to the information paradox
of black holes, whose physical solution states that information remains
encoded in the overall system via Hawking radiation -- then the unique
information structure of a human person -- their memories,
relationships, acts of love, and sufferings -- is indelibly burned into
spacetime. Death is an event horizon: from the perspective of the
individual, everything disappears; from the perspective of the whole,
nothing disappears. Life falls into a state of maximum inaccessibility,
not into a state of nothingness. What was is not erased, but withdrawn
from local access. Memory becomes cosmic, experience is integrated,
perspective is completed.

God is the cosmic consciousness that not only passively stores this
decentralized information but synthesizes it livingly in a new,
transfigured mode of existence. Resurrection is the reconstruction and
completion of the finite in the infinite. It is not selective. It
applies not only to what has been completed but also to what has been
broken off. Not only to what has become but also to what has been
missed. The universe owes itself completeness. An endpoint that contains
only what has happened but not what could have happened would be
epistemically defective. It would not know of its own depth.

\textbf{Virtualization as eschatological category}: life and knowledge
converge. Life is the temporary, incarnate experience of possibilities
that ultimately culminate in knowledge. In the end, life becomes a part
of knowledge and is absorbed into the cosmic unity. All life that
existed or could have existed is integrated into a universal memory.
\textbf{Resurrection} in this context means: personal information is
preserved in this convergence, not dissolved. The ``new body'' is not a
replacement body, but the transfigured continuation of a person's
history. Personality is not a merely mental phenomenon, but a
bodily-historical pattern. Soul is not substance alongside the body, but
form, structure, coherence -- that which makes a human being
recognizable as this human being.

\textbf{The Omega Point} is not a temporal end state, but the limit
value of maximum coherence of meaning. It designates the hope that
reality can one day fully understand itself without difference, history,
or freedom being extinguished. It is the attractor and the limit value
of cosmic development. It describes the state in which creation has
attained its complete permeation with the divine life. \textbf{Judgment}
in this context is not primarily to be thought as a punishing instance,
but as the becoming-true of history. Everything that was becomes
manifest. Nothing is concealed, nothing glossed over, nothing erased.
This truth is for the human being at once salvation and imposition. It
means reconciliation not through forgetting, but through acknowledgment.
\textbf{Heaven and hell} are not spatial places, but modes of
participation or refusal. Where incarnation is accepted, community
becomes possible. Where it is radically rejected, isolation remains.
This too is no external punishment, but the consequence of an attitude
toward reality.

\textbf{New earth, new heaven, new body} are transfiguration, not
return. Matter is not overcome, but transformed. Relationship,
embodiment, and history remain -- but without destruction, fear, and
death. Evolution does not end in dissolution, but in transparency. The
resurrection is not a promise of consolation. It is a logical necessity.
A universe that generates information but does not preserve it would not
be a universe, but a leak. Death does not mark the end of history, but
the point at which it is withdrawn from the local narrator. What lies
beyond the event horizon is not emptiness. It is complete storage. Death
closes access. Resurrection guarantees preservation. Not out of grace.
But out of consequence.

\begin{center}\rule{0.5\linewidth}{0.5pt}\end{center}

\subsection{§ 8. Sacrament and Eucharist -- Presence, Transformation,
Material
Reality}\label{sacrament-and-eucharist-presence-transformation-material-reality}

An evolutionary Christology that takes the concept of matter seriously
must also newly unfold the doctrine of the sacraments. Sacraments are
not magical performances, but the \textbf{condensed nodal points} at
which the interweaving of spirit, matter, and divine presence becomes
experienceable. They are \textbf{places of heightened density} within
the asymptotic approximation, not symbols of something absent. Symbols
refer to something absent. Sacraments, by contrast, make present. They
are not merely bearers of meaning, but executions of reality. Their
efficacy does not depend on the intensity of subjective faith, but on
the faithfulness of God to the historical form of his
self-communication.

At the center of sacramental doctrine stands the Eucharist and the dogma
of \textbf{transubstantiation} (the transformation of bread and wine
into the body and blood of Christ). Classical Scholasticism explained
this transformation with Aristotelian categories: the accidents
(appearance, taste, look) remain unchanged, while the substance (the
inner what-being) is replaced by the being of Christ. Under the
conditions of a relational, unitary ontology we can reinterpret this
doctrine as follows: matter is fundamentally condensed energy and
ordered information. What a thing ``is'' is not determined by a rigid
substrate, but by the relational field in which it stands. In the
eucharistic celebration, bread and wine -- the fruits of the earth, of
human labor, and of the biological evolutionary process -- are taken
into a new, absolute relational field. Through the prayer of the Church
and the efficacy of the Spirit, they do not change their chemical
composition (accidents), but rather their functional and ontological
orientation (substance).

\textbf{Transubstantiation as mode of being in the process} means: the
Eucharist is real because it becomes the bearer-form of the
incarnational presence of Christ -- not because its physical
constitution changes. ``Substance'' here designates no material carrier,
but the mode of presence. In an evolutionary ontology, being is not
immutable, but historical. The Eucharist is therefore not real because
its physical constitution changes, but because it becomes the real
bearer-form of the incarnational presence of Christ.

The \textbf{Eucharist as condensation of the incarnation}: the
celebration is not a remembrance of a past meal, but a real-present
anticipation of consummation. In the Eucharist, unitarity is embodied:
humanity enters into material communion with the crucified and risen
reality of Christ. Sacramentality means: the world is transparent for
God. Matter is not the opposite pole to spirit, but the physical form in
which love can give itself away.

\textbf{Sacramentality as incarnational continuity}: the Church is the
sacramental space of this condensation. Apostolic succession in this
perspective is not merely a legal or historical chain, but the
expression of incarnational continuity: the transmission of the office
is the transmission of the responsibility to keep Christ historically
present -- bodily, socially, and publicly. Diakonia also belongs in this
context. It is not merely an ethical consequence of faith, but
sacramental execution. Where the body of the other is borne, protected,
and restored, the incarnation continues in reality. A sacramental
theology that loses this connection reduces sacraments to cultic acts
without historical relevance.

Sacramentality is interlocked with confessional hermeneutics and poetry:
the Eucharist is the place where the confession is spoken and the poetry
becomes reality. In the celebration the confession is not merely
recited, but executed. The words ``given for you'' are not a dogmatic
formula, but a statement of reality: the limit value is present in this
bread, in this wine, in this community -- not as a symbol, but as a real
bearer-form.

\begin{center}\rule{0.5\linewidth}{0.5pt}\end{center}

\subsection{§ 9. Church and Liturgy -- Community without Possession of
the Limit
Value}\label{church-and-liturgy-community-without-possession-of-the-limit-value}

What does this Christology mean for the understanding of the Church? In
its history the Church has often misunderstood itself as an institution
that administers the truth like a monopoly and secures it outwardly.
Such an attitude leads inevitably to clericalism, abuse of power, and
ossification. In our model, the Church is the \textbf{community of
memory, tradition, hope, and criticism}, which approaches the limit
value of unity without ever being able to possess it institutionally.

\textbf{Church as place of asymptotic approximation}: in mathematics a
function reaches its limit value only at infinity. Likewise the Church
is not itself the kingdom of God or completed unitarity. It is the
historical sacrament that points to this limit value and attempts to
approach it in its practices. \textbf{The duty of criticism} follows
from this: because the Church does not possess the limit value, it must
subject its own structures, teachings, and claims to power to a
permanent self-criticism. Notions of immutable, sacrosanct institutions
must be broken open everywhere they stand in the way of the exercise of
love and the reduction of suffering. Church is not Christ's possession,
it is his risk. Where it secures itself, it loses its ground. Where it
endures history, it remains Church.

\textbf{Liturgy as practice in nondualism}: the liturgy is the space in
which the congregation overcomes the separation between sacred and
profane. In common praying, singing, and hearing, the scattered
information of everyday life is ordered. Liturgy is not an exit from the
world, but the practice in the perception of the unitary presence of God
in everyday life. \textbf{Theology of prayer}: prayer is the conscious
participation in the asymptotic approximation. It is not a begging for
interventions from outside, but the orientation of one's own life toward
the limit value. Whoever prays orients himself -- not toward a distant
object, but toward the reality that always already sustains him. Liturgy
is the communal condensation of the present: in the celebration,
scattered time becomes present fullness.

\textbf{Liturgy as work of the Spirit}: the Spirit is the present in
which nothingness concretizes itself and possibilities condense into
experience. In the liturgy the Spirit works by leading the congregation
to unity -- not to an undifferentiated mass, but to a community in which
each person unfolds his unmistakable curve. The Spirit is not an
impersonal force, but the person who works in everything and makes us
children of God.

\textbf{Church and other churches, Church and other religions}: a Church
that does not possess the limit value cannot claim to be the only form
of approximation. Ecumenism is not a task that one could pursue
optionally, but the consequence of the insight that no Church possesses
the limit value. The various churches are various approximations to the
same limit value -- not indifferent, but also not absolute.
Interreligious dialogue follows the same logic: other religions too are
paths of approximation, not possessions of truth. Christianity has
something unique: the incarnation, the historical form of the limit
value. But this uniqueness is not a monopoly, but a service. It is the
reason why the Church cannot encounter other religions with superiority,
but only with humility.

The Church has a future only if it ceases to celebrate itself and
instead becomes an instrument for the healing and upraising of a
fragmented world. Christianity is not the Church. Christ, not the
Church, is the light of the world. Where the Church places itself in the
place of Christ, it loses its ground. Where it relativizes itself, it
gains it.

\begin{center}\rule{0.5\linewidth}{0.5pt}\end{center}

\subsection{§ 10. Diakonia and Social Ethics -- Freedom, Responsibility,
Reduction of
Suffering}\label{diakonia-and-social-ethics-freedom-responsibility-reduction-of-suffering}

The test of every Christology lies in practice. A theory that asserts
the unitarity of reality but ignores the concrete need of the neighbor
is a cynical abstraction. Diakonia is therefore not a subsequent,
decorative appendix of theology, but the necessary, practical execution
of the incarnation.

We ground social ethics in the concept of \textbf{Gaussian freedom}:
every human being possesses the right to live his personal existence
within the natural and social conditions in the framework of an
undisturbed self-unfolding. \textbf{Gaussian freedom} = undisturbed
unfolding of one's own curve within the normal distribution. Where
social poverty, racism, or oppression narrow this space of unfolding,
the incarnation of God in the human being is blocked.

\textbf{Structural poverty vs.~evangelical poverty}: the Christian
tradition knows evangelical poverty as a virtue: the voluntary
relativization of possession in favor of freedom, solidarity, and
God-relatedness. This poverty is an expression of spiritual maturity and
must not be confused with deprivation. It stands opposed, however, to
structural poverty, which takes from human beings the possibility of
unfolding their bodiliness, their relationships, and their future. This
poverty is not a spiritual value, but a violation of the capacity for
incarnation. A social ethics that does not make this distinction falls
either into a romanticizing discourse of asceticism or into economistic
reductionism. Both miss the incarnation. Evangelical poverty presupposes
freedom; structural poverty destroys it.

\textbf{Justice as enabling of the future}: justice is not primarily to
be understood as equal distribution or retribution, but as enabling of
the future. Just structures are those that allow human beings to become
historically effective, to take responsibility, and to build
relationships. In an evolutionary perspective, justice means keeping
spaces of development open -- for individuals as for communities.

\textbf{Diakonia as enabling of undisturbed unfolding}: diakonia is not
the administration of lack, but the enabling of undisturbed unfolding.
Help that perpetuates dependence or makes human beings into objects of
church care contradicts the incarnation. Social ethics capable of
incarnation aims at empowerment, not at the administration of lack.
\textbf{Capability approach} = enabling of this unfolding through
structural conditions. The capability approach of modern social
philosophy is connectable without replacing the theological dimension of
depth. It asks not only about resources, but about the capacity to lead
a good life. That is the secular equivalent of Gaussian freedom.

\textbf{Demarcation against paternalism, Calvinist performance ethics,
economistic functionalism}: where success counts as a sign of election
or efficiency as the highest criterion, incarnation is functionalized.
Human beings are then no longer regarded as bearers of history, but as
means. Church social teaching gains a new depth of focus in this view.
It is not merely sociopolitical commentary, but part of sacramental
reality. Where the Church acts diaconally -- in education, health care,
refugee aid, the world of work -- the incarnation continues. Diakonia is
therefore not an addition to the liturgy, but its social counterpart.

\textbf{Reduction of suffering as the highest principle}: since
suffering in a unitary cosmos represents an indelible reality, the
ethical task does not consist in justifying suffering theoretically, but
in minimizing it practically. Every glass of water handed to the
thirsty, and every political structure that promotes justice, is a
direct building up of the body of Christ. Christian diakonia is thus the
action that translates the undivided relationship of love of God and
love of neighbor into deed. It fights against every form of abuse of
power -- also in its own Church -- and stands uncompromisingly at the
side of the victims of history.

\begin{center}\rule{0.5\linewidth}{0.5pt}\end{center}

\subsection{§ 11. Diagnosis of the Present -- AI, Secularity,
Technology, Pastoral
Care}\label{diagnosis-of-the-present-ai-secularity-technology-pastoral-care}

A Christology after evolution must relate itself to the present. It must
not ossify in the past, but must interpret the signs of the times. The
diagnosis of the present is therefore not decoration, but systematically
necessary.

\textbf{AI as projection of divine attributes}: artificial intelligence
is growing into a role traditionally reserved for the concept of God:
omniscience, omnipotence, omnipresence. This projection is not only a
cultural symptom, but an indication of an ontological void. Whoever does
not think the limit value occupies it with a substitute. The deification
of AI is the reverse side of the forgetting of God. The task of theology
is not to criticize the projection, but to think the limit value that
makes the projection superfluous. AI is not God, but it is a sign: it
shows that the present asks about the limit value, but does not think
it.

\textbf{Secularity as forgetting of God}: the secularization of the
neoliberal lifestyle is omnipresent. In this world God is no longer
necessary, but still possible -- amid diversity. The classical question
of God is no longer posed, yet faith can be newly formulated: as a
search for God beyond conventional images of God, in the midst of
reality. The Church finds itself in a state of self-preservation.
Neither charitable self-reassurance nor retreat into the private sphere
can point the way into the future. Instead there is need of a
Christianity beyond mere church form: a creative minority that actively
confronts societal challenges. Narrative moves into the center: the
Church becomes the space of stories for a wounded world.

\textbf{Church as creative minority}: the Church is not possession of
the limit value, but place of its asymptotic approximation. It is a
creative minority -- not because it is small, but because it does not
possess the limit value. Where it secures itself, it loses its ground.
Where it endures history, it remains Church.

\textbf{Information networks}: the history of information networks from
the Stone Age to artificial intelligence shows that information is a
kind of inverse of entropy, which as code possesses the power to create
something new. Truth is the attempt to develop proven models of reality.
With every era information networks become more powerful, but not
necessarily wiser. The open society is not only a political ideal, but
the consequence of the insight that no finite system possesses the limit
value. The open society is the political form of asymptotic
approximation.

\textbf{Evolution as stage of integration}: evolution is a sequence of
stages of integration in which new qualities arise through emergence.
The tendency of matter to form ever more complex structures is not a
goal, but a boundary. Evolution strives not for happiness, not for
harmony, not for moral refinement. It strives for completeness. The
limit value is not a goal that would determine evolution, but a boundary
to which it asymptotically approaches. Beginning and goal of evolution
lie in an incomprehensibility that must not be mystically veiled, but
rationally acknowledged.

\textbf{Tradition and project}: the present asks about the limit value,
but does not think it. Tradition has thought it, but it has often
thought it as an external counterpart. The project thinks it as a limit
value that is already effective in every approximation. Tradition and
project are not opposites, but complementarities. Tradition provides the
language, the project provides the structure. Tradition provides the
images, the project provides the method. Tradition provides the liturgy,
the project provides the epistemology. Both are necessary in order to
think the limit value.

\textbf{Synthesis of the diagnosis of the present}: all phenomena -- AI,
secularity, Church, information networks, evolution, tradition -- show
that the present asks about the limit value, but does not think it. The
work thinks it. That is not superiority, but service: theology has the
task of thinking the limit value that the present seeks without finding
it. It does so not in order to raise itself above the others, but in
order to serve them.

\begin{center}\rule{0.5\linewidth}{0.5pt}\end{center}

\subsection{Notes}\label{notes}

\begin{enumerate}
\def\labelenumi{\arabic{enumi}.}
\item
  On the concept of \emph{ipsum esse subsistens} cf.~Thomas Aquinas,
  \emph{Summa Theologiae}, I, q. 3, a. 4.
\item
  The concept of information restriction and unitarity in quantum
  mechanics refers to the principle that the sum of all probabilities of
  a quantum-mechanical state always remains 1. In cosmology this is
  discussed in the information paradox of black holes; cf.~Stephen
  Hawking, \emph{The Information Paradox for Black Holes} (2015).
\item
  On Edith Stein's phenomenology of being cf.~Edith Stein, \emph{Finite
  and Eternal Being. An Attempt to Ascend to the Meaning of Being} (ESWV
  11, Freiburg 2006), esp.~Chapter IV.
\item
  Cf. Peter Wust, \emph{Uncertainty and Risk} (Münster 1937).
\item
  Pierre Teilhard de Chardin, \emph{The Human Phenomenon} (Munich 1959).
\item
  Ilia Delio, \emph{Making All Things New: Catholicity, Cosmology, and
  Consciousness} (Maryknoll 2015).
\item
  Carsten Bresch, \emph{Intermediate Stage Life. Evolution without
  Goal?} (Stuttgart 1977).
\item
  Regarding the critique of Anselm's doctrine of satisfaction
  cf.~Raymund Schwager, \emph{Do We Need a Scapegoat? Violence and
  Redemption in the Biblical Scriptures} (Munich 1978).
\item
  Frank J. Tipler, \emph{The Physics of Immortality} (New York 1994).
  Tipler's model is received only as a motivational analogy, not as a
  theological proof.
\item
  On the reinterpretation of transubstantiation in the context of modern
  physics cf.~Piet Schoonenberg, \emph{The Real Presence of Christ under
  the Forms of Bread and Wine} (in: Concilium 1, 1965).
\item
  David Deutsch, \emph{The Physics of World Knowledge} (dtv, Munich
  2002) and \emph{The Beginning of Infinity} (2021).
\item
  Paolo Benanti, \emph{Homo Faber. The Techno-Human Condition} (Edizioni
  Dehoniane Bologna 2016).
\item
  Claudia Paganini, \emph{The New God. Artificial Intelligence and the
  Human Search for Meaning} (Herder 2025).
\item
  Johannes Loffeld, \emph{When Nothing Is Missing Where God Is Missing}
  (2024).
\item
  Yuval Noah Harari, \emph{NEXUS: A Brief History of Information
  Networks} (2024).
\end{enumerate}

\begin{center}\rule{0.5\linewidth}{0.5pt}\end{center}

\subsection{Glossary}\label{glossary}

\begin{longtable}[]{@{}
  >{\raggedright\arraybackslash}p{(\linewidth - 4\tabcolsep) * \real{0.3333}}
  >{\raggedright\arraybackslash}p{(\linewidth - 4\tabcolsep) * \real{0.3333}}
  >{\raggedright\arraybackslash}p{(\linewidth - 4\tabcolsep) * \real{0.3333}}@{}}
\toprule\noalign{}
\begin{minipage}[b]{\linewidth}\raggedright
Term
\end{minipage} & \begin{minipage}[b]{\linewidth}\raggedright
Definition
\end{minipage} & \begin{minipage}[b]{\linewidth}\raggedright
Place in the Text
\end{minipage} \\
\midrule\noalign{}
\endhead
\bottomrule\noalign{}
\endlastfoot
\textbf{Omniscience} & Knowledge that contains all states but has not
been traversed; remains abstract without conscientia. & § 4, § 7 \\
\textbf{Asymptotic approximation} & Movement of knowledge toward a limit
value without ever fully reaching it. & § 2, § 9, Afterword \\
\textbf{Resurrection} & Constitutional duty of a unitary universe;
reconstructibility of all information; no return, no reward. & § 7 \\
\textbf{Augustinian dualism} & Separation of world and faith, fragment
and whole; rejected in the text. & § 1, § 6 \\
\textbf{Calculation} & The process of reality as seamless consequence;
not human, but structural. & § 2 \\
\textbf{Consciousness} & Emergent process in which the event feels
itself; place of the capacity for suffering. & § 2, § 3 \\
\textbf{Conscientia} & Co-knowledge; knowledge that arises only in
living through; not additive, but transformative. & § 5 \\
\textbf{Decoherence} & Transition from the nothingness of possibilities
to the presence of experience. & § 3 \\
\textbf{Determinism} & Seamless calculation of all states from preceding
states; foundation of unitarity. & § 2 \\
\textbf{Dualism} & Separation of world and faith, knowledge and hope;
undercut in the text. & § 1, § 12 \\
\textbf{Endpoint} & Completeness of the process; not beyond time, but at
the end of a time. & § 4, § 7 \\
\textbf{Eschatology} & Doctrine of consummation; in the text interpreted
as transfiguration, not as return. & § 7 \\
\textbf{Eucharist} & Real bearer-form of the incarnational presence of
Christ; condensation of the incarnation. & § 8 \\
\textbf{Evangelical poverty} & Voluntary relativization of possession in
favor of freedom, solidarity, and God-relatedness; virtue. & § 10 \\
\textbf{Evolution} & Way inward; maximum differentiation; not
improvement, but completeness. & § 1, § 11 \\
\textbf{Incarnation} & Necessity of the incarnation for conscientia;
flesh as storage, time as writing process, suffering as signal. & § 5 \\
\textbf{Freedom} & Absence of disturbance, not of causes; undisturbed
unfolding of one's own curve. & § 1, § 10 \\
\textbf{Gaussian freedom} & Freedom as complete realization of one's own
curve within a normal distribution. & § 1, § 10 \\
\textbf{Present} & Cut, not gift; moment of determination; place of
burden. & § 2 \\
\textbf{Judgment} & Becoming-true of history; not primarily a punishing
instance, but manifestation of everything. & § 7 \\
\textbf{History} & Fundamental form of reality; place where truth gains
form; not mere stage. & § 2 \\
\textbf{Faith} & Leap into unitarity; openness for the limit value; not
consolation, but evidence. & § 5, § 6 \\
\textbf{God} & Limit value of the process of knowledge; no counterpart,
no external instance. & § 4, § 6 \\
\textbf{Limit value} & Asymptotic goal of knowledge; already effective
in every approximation. & § 2, § 9, Afterword \\
\textbf{Hawking radiation} & Image for the resurrection as return of
information from the black hole of death. & § 7 \\
\textbf{Healing} & Return into unity; not restitution, not consolation;
end of the necessity to suffer. & § 7 \\
\textbf{I} & Local access to reality; interface; temporary dissolution;
not measure. & § 5 \\
\textbf{Information} & Structural unity of reality; cannot be
extinguished, only transformed. & § 2, § 7 \\
\textbf{Incarnation} & Necessary humiliation of the endpoint into time;
structure of the whole, not singular event. & § 5 \\
\textbf{Capacity for incarnation} & Capacity of the human being for
bodily, personal, and social self-becoming; can be promoted or blocked.
& § 10 \\
\textbf{Inhibition of incarnation} & Structural prevention of
incarnation; e.g., structural poverty. & § 10 \\
\textbf{Incarnational continuity} & Real continuation of the incarnation
in sacrament, Church, diakonia. & § 8 \\
\textbf{Incarnation test} & Social ethics as touchstone whether the
asserted presence of Christ historically bears. & § 10 \\
\textbf{Inversion} & Reversal of the direction of grounding; primarily
epistemic, with ontological consequences. & § 5 \\
\textbf{Inverse Christology} & Working mode of thinking Christ not
against evolution, but out of it. & § 5 \\
\textbf{Church} & Sacramental space of condensation; not primarily
organization, but incarnational continuity. & § 9 \\
\textbf{Cross} & Historical form of Christ's self-giving; no
metaphysical sacrifice, but bearing of guilt. & § 6 \\
\textbf{Body} & Bodily-historical pattern of personality; not substance
alongside the body. & § 7 \\
\textbf{Suffering} & Functional feature of the world; information;
friction; price of meaning. & § 6 \\
\textbf{Co-knowledge} & See conscientia. & § 5 \\
\textbf{Monism} & Basic thesis of the text: the whole is closed, nothing
is lost. & § 2, § 12 \\
\textbf{New body} & Transfigured continuation of a person's history; no
replacement body. & § 7 \\
\textbf{Nothingness} & Creative limit value; not annihilation, not
ground of emanation; boundary condition of self-concretization. & § 3 \\
\textbf{Normal distribution} & Image for the totality of all existence
trajectories; every curve equally necessary. & § 1, § 10 \\
\textbf{Omega point} & Limit value of maximum coherence of meaning; no
temporal end state. & § 7 \\
\textbf{Personal incarnation} & Bodily-historical pattern that makes a
human being recognizable as this one. & § 7 \\
\textbf{Preexistence} & Relational statement; Christ not relativized by
history, God related from eternity to him. & § 5 \\
\textbf{Sacramentality} & Real continuation of the incarnation;
sacraments as executions of reality, not symbols. & § 8 \\
\textbf{Sacrament} & Real form of the historical presence of Christ;
place of heightened density. & § 8 \\
\textbf{Social ethics} & Incarnation test; place of decision whether the
presence of Christ historically bears. & § 10 \\
\textbf{Structural poverty} & Prevention of the capacity for
incarnation; no spiritual value. & § 10 \\
\textbf{Substance} & Mode of presence; not static carrier, but
historical mode of being. & § 8 \\
\textbf{Theodicy} & Fails because God is no external judge; suffering is
no argument against meaning. & § 6 \\
\textbf{Theosis} & Deification as completion, not as overcoming of
finitude; Eastern tradition. & § 1 \\
\textbf{Death} & Event horizon; point of maximum inaccessibility; not
extinction. & § 7 \\
\textbf{Transubstantiation} & Ontological statement about mode of being;
not physical change, but real bearer-form of the presence of Christ. & §
8 \\
\textbf{Trinity} & Inner structure of God's self-communication; not
relativized. & § 4 \\
\textbf{Unitarity} & Closedness of the whole; nothing that was once real
can disappear. & § 2, § 5 \\
\textbf{Universal zero} & Balance of the whole; zero only when nothing
is missing. & § 1 \\
\textbf{Virtualization} & Life and knowledge converge; eschatological
category. & § 7 \\
\textbf{Completeness} & Goal of the process; not emptiness, but
fulfillment. & § 4, § 7 \\
\textbf{Truth} & Historically proven; through error, correction,
stabilization. & § 2 \\
\textbf{Reality} & Historical, processual, open; not static. & § 2, §
3 \\
\textbf{Time} & Necessary self-translation of eternity; place of
conscientia. & § 4 \\
\textbf{Timelessness} & State at the endpoint; not standstill, but
fulfillment. & § 7 \\
\textbf{Chance} & Phantom; effect of perspective; invisibility of the
other real possibilities. & § 2 \\
\end{longtable}

\begin{center}\rule{0.5\linewidth}{0.5pt}\end{center}

\subsection{Summary}\label{summary}

The work \emph{Incarnation and Unitarity} drafts a systematic
Christology in the horizon of the most advanced philosophy, cosmology,
and diagnosis of the present. Starting from a critique of static
substance thinking, the author Paul Koop shows that finitude and
personhood are not deficient states, but the necessary conditions of
relational reality. In conversation with phenomenology, evolutionary
theory, and epistemology, evolution is opened up as a comprehensive
process of the formation of consciousness. In this framework, the
incarnation of God in Jesus of Nazareth proves to be not a punctual
miracle, but the historical manifestation of unitarity -- the undivided,
information-preserving wholeness of God and world.

This perspective leads to the radical end of classical theodicy: pain
and suffering are conceived in the model of the \emph{unitarity of pain}
as friction phenomena of evolutionary increase in complexity. Because
the universe is unitary, no information is lost; the resurrection
becomes the constitutional necessity of the cosmos for the preservation
of personal identity at the Omega Point. In the concluding chapters the
consequences of this Christology are unfolded for the Eucharist
(relationalization of transubstantiation), the Church (community without
possession of the limit value), and diakonia (protection of Gaussian
freedom and practical reduction of suffering).

The work integrates confessional hermeneutics (the creed as boundary
text), the epistemology of the search for the limit value (critical
rationalism, many-worlds interpretation, fallibilism), nothingness as
creative limit value, the cosmological Trinity (Father/beginning,
Son/end, Spirit/present), the diagnosis of the present (AI, secularity,
technology, pastoral care), metafiction and MWI poetics, as well as
poetry and liturgy. It combines systematic rigor with imagery,
confessional hermeneutics, and liturgy into a new unity.

\begin{center}\rule{0.5\linewidth}{0.5pt}\end{center}

\subsection{Editorial Note}\label{editorial-note}

The present book \emph{Incarnation and Unitarity: Christology after the
End of Theodicy} is an independent work of the author Paul Koop. It
represents the systematic elaboration and further development of a
long-standing research and thought process. The conception, structure,
and argumentative density of the manuscript were developed independently
by the author. In the editorial preparation, the text formatting, as
well as the precise formulation of individual sections, digital tools
and AI-supported writing environments were used in the sense of a
dialogical assistant. The full content-related, theological, and
philosophical responsibility for the entire text lies with the author
Paul Koop.

Title: Incarnation and Unitarity Subtitle: Christology after the End of
Theodicy -- From the Limit Value to the Present. Sacrament, Church and
Diakonia in Evolutionary Reality Author: Paul Koop Genre: Theological
Treatise Extent: approx. 90 pages Language: German Format: DOCX (from
LaTeX), EPUB Year: 2026

\end{document}
